%
% atomic-page-write.spec.tex
%
% Atomic page write: a Z specification of how ChunkStoreFunnel writes,
% replaces, and cleans up one document's chunks, beside the directory-sync
% writer and its registry (quarry-70zm).
% Companion to docs/design/atomic-page-write.md, which cites the sections
% below by number and records the tool runs.
%
% Check:        z-spec check docs/design/atomic-page-write.spec.tex
% Model-check:  probcli with -p MAX_INITIALISATIONS 100 (the command lines
%               are in atomic-page-write.spec.report.json)
%
% Must be .tex, not .md: the z-spec tools glob for .tex.
%

\documentclass[a4paper,10pt,fleqn]{article}
\usepackage[margin=1in]{geometry}
\usepackage{fuzz}
\usepackage{booktabs}

\newcommand{\field}[1]{\texttt{#1}}

\begin{document}

\title{Atomic Page Write: A Z Specification}
\author{Guido R, for Punt Labs}
\date{October 2026}
\maketitle

% ===================================================================
\section{Scope and modelling choices}
\label{sec:scope}
% ===================================================================

This specification models the rows of \emph{one document name} in the
\field{chunks} table and the sync registry's claim on that name. Four kinds of
actor touch them:

\begin{itemize}
  \item up to two \emph{funnel writers} (\field{wa}, \field{wb}) that replace,
    create, or heal the name, and that may fail or crash at any step;
  \item the \emph{directory-sync writer}, which owns the names its registered
    directories cover, writes its rows window by window, and deletes them by
    name;
  \item the \emph{operator}, who deregisters a directory with its data purged
    (\field{Purge}); registering a directory claims no name until sync writes
    that file (the registry lists the name only then), and deregistering with
    the data kept leaves the name in the retained set, which is ownership that
    does not change, so neither is an operation;
  \item the \emph{environment}, which may delete a single row of a funnel
    generation behind everyone's back, producing a damaged generation.
\end{itemize}

The rows of every other document are irrelevant to the properties checked,
because every operation below is predicated on one name; that is itself
property I4. \textbf{What the model does not cover} is listed in
\S\ref{sec:nongoals}; every result in the design is scoped to this model.

\subsection{Generations, keys, and the five row classes}

\field{rows}~$g$ is how many of the $K=2$ chunks of class $g$ are in the table.

\begin{center}
\begin{tabular}{@{}cl@{}}
\toprule
$g$ & Rows of that class \\
\midrule
0 & funnel rows with no key (written before this protocol): \emph{legacy}; or,
    when \field{g0sync} holds, the \emph{unmarked sync rows} an earlier build
    wrote \\
1, 2 & funnel generations with an ordinary write key ($1<2$ in capture order) \\
3 & a funnel generation whose key lies beyond \emph{twice} the skew
    allowance: the \emph{future} key (a clock stepped back, a corrupt row) \\
4 & sync-marked rows (\field{write\_id = 'sync'}) \\
\bottomrule
\end{tabular}
\end{center}

A funnel generation $g\in 1\upto2$ is \emph{complete} when
\field{rows}~$g=K$. Key~2 doubles as the \emph{write ceiling}: the highest key
a write may carry (writes clamp at $\mathit{now}+\mathit{allowance}$), while a
stored key is accepted up to twice the allowance, so a stored copy can sit at or
above the ceiling and a submission write cannot exceed it (\field{RefuseCeiling}).
Generation~3 is \emph{never} complete: a key beyond \emph{twice} the allowance is
damaged by classification, contributes no stamp, is never held, and is
reapable once a complete generation exists. New writers never choose key~3
(design \S7), so the model gives writers keys $1\upto2$ only.

A funnel writer picks its key at \emph{begin}, \emph{independently of begin
order}. Because the implementation adds a generation with a single LanceDB
commit, a protocol step can only produce $0$ or $K$ rows for $g\in1\upto2$;
the model does not assume that, and records whether a protocol step
\emph{produces} a partial generation (\field{partialVisible}).

\subsection{Ownership}

The sync registry's claim on the name is \field{reg}: the \emph{files} table
lists it or the retained set does. It becomes true when sync first writes the
name. A name is \emph{owned} when \field{reg} holds or any marked sync row
exists. If funnel rows exist when sync first writes the name (the order
reversal), sync \emph{takes the name over} --- deletes the funnel rows by name
--- only when those rows are \emph{file-sourced from the same path}
(\field{fsrc}); otherwise (a label, a capture, a URL page) sync skips the file
and records a conflict (\field{SyncConflict}), and the funnel copy stays. A
funnel write, and every heal (the writer's, a held-check's, the sweep's), begin
only when the name is not owned; if the registry lookup fails the writer fails
closed (\field{LookupFails}).

\subsection{Two protocols, one switch}

\field{mode}~$=0$ is the protocol proposed in the design. \field{mode}~$=1$ is
the funnel as it was (overwrite-delete by name, windowed insert into the real
name, cleanup by name, create-if-absent by name existence, no ordering, no
ownership). The same witnesses are checked against both; mode~1 must
\emph{fail} them and mode~0 must not.

\subsection{Re-verification before the delete}

An earlier run of this model found a counter-example: a healer observes key~2
complete, the environment removes one row of key~2, and the healer's delete
(older and damaged generations) then removes the older complete copy too. The
delete therefore \emph{re-verifies} the observed generation immediately before
deleting and does nothing if it is no longer complete. Re-verify and delete are
modelled as one step; in the implementation they are two calls, so a second
external fault inside that millisecond is a residual double fault.

\subsection{Witnesses, not invariants}

A schema invariant that an operation can violate makes that operation
unsatisfiable under Z's schema calculus: ProB would then report nothing
rather than a violation. Each property is therefore a latched witness
\field{ZBOOL} that an operation sets when it commits the bad transition,
checked by asking ProB whether any state with the witness true is reachable.

\subsection{Non-goals: what a model result does not claim}
\label{sec:nongoals}

\begin{itemize}
  \item Row content, ordering inside a generation, and memory (I5).
  \item Legacy rows that are truncated but look complete.
  \item The comparison of a new key against a legacy copy's
    \field{ingestion\_timestamp}; the model orders generation~0 below every
    real key. The skew allowance is a fixed boundary: keys are either ordinary
    ($1\upto2$) or future ($3$); skew \emph{within} the allowance, which can
    misorder two writes, is not modelled.
  \item The registry lookup and the commit are not one atomic step across
    processes. \field{Add}'s guard reads the registry and the table in one
    step; the implementation reads the registry inside the in-process write
    lock, and a registration that lands after the read and before the commit
    is a residual (sync then takes the name over at its first write).
  \item Mixed-version writers beyond the environment's single-row delete.
  \item That the heal's delete predicate is the one the model writes. The model
    fixes the predicate (never rows of class 4, never unmarked sync rows, never
    when the name is owned at heal start); the implementation must match it and
    the design's tests pin that.
  \item That the sync operation is as strong in code as here. \field{SyncInsert}
    decides the takeover and writes in one step; the code checks, deletes, and
    inserts as separate calls, and a crash between them leaves a partial
    sync-owned document, which the registry watermark resumes.
  \item That the registry lookup \emph{fails closed}. \field{LookupFails} is the
    only transition available when the lookup fails, and it changes no row; the
    model cannot show that an implementation does not fail open.
  \item The sync path's own protocol (registry watermark, tail delete).
  \item Liveness beyond a \emph{possibility} claim: the CTL property below says a
    clean quiescent state is reachable from every reachable state, not that
    every execution gets there.
\end{itemize}

% ===================================================================
\section{Types and constants}
\label{sec:types}
% ===================================================================

\begin{zed}
ZBOOL ::= ztrue | zfalse
\end{zed}

\begin{zed}
WRITER ::= wa | wb
\end{zed}

A writer's phase. \field{staging}: embedding windows into a private spool,
invisible to readers (the windows are not modelled: the spool is invisible, so
it is one phase). \field{ready}: spool complete, nothing committed.
\field{healing}: about to heal. \field{observed}: the newest complete generation
has been read, the delete not yet issued. \field{inserting}: the old protocol,
writing into the real name window by window.

\begin{zed}
PC ::= idle | staging | ready | healing | observed | inserting
     | done | failed | crashed
\end{zed}

\begin{axdef}
mode : \nat
\where
mode = 0
\end{axdef}

% ===================================================================
\section{State}
\label{sec:state}
% ===================================================================

\begin{schema}{Store}
rows : (0 \upto 4) \fun (0 \upto 2) \\
reg, g0sync, fsrc : ZBOOL \\
pc : WRITER \fun PC \\
wgen : WRITER \fun (1 \upto 2) \\
seen : WRITER \fun (0 \upto 5) \\
partialVisible, lost, unsafeDelete, skippedOnPartial, inversion : ZBOOL \\
syncTouched, syncClash, labelLost : ZBOOL
\end{schema}

\begin{itemize}
  \item \field{rows}, \field{reg}, \field{g0sync}: as above.
  \item \field{fsrc}: any funnel rows now present are file-sourced from the same
    source path as sync's file (meaningful only when funnel rows exist).
  \item \field{wgen}: the write key a funnel writer chose at begin.
  \item \field{seen}: the newest complete generation a healer observed; $5$
    means \emph{none was complete}.
  \item Witnesses (latched, never reset): \field{partialVisible} (I1),
    \field{lost} (I2), \field{unsafeDelete} (I4), \field{skippedOnPartial}
    (create-if-absent), \field{inversion} (I6), \field{syncTouched} (a funnel
    or heal step changed marked sync rows, or unmarked sync rows), and
    \field{syncClash} (a funnel step added a generation to an owned name), and
    \field{labelLost} (sync's takeover deleted funnel rows that were not
    file-sourced from the same path: a label, capture, or URL page).
\end{itemize}

\begin{schema}{Init}
Store'
\where
rows'~0 \in 0 \upto 2 \\
rows'~1 = 0 \\
rows'~2 = 0 \\
rows'~3 \in \{ 0, 2 \} \\
rows'~4 \in 0 \upto 2 \\
g0sync' = ztrue \implies reg' = ztrue \\
fsrc' \in ZBOOL \\
\forall v : WRITER @ pc'~v = idle \land wgen'~v = 1 \land seen'~v = 5 \\
partialVisible' = zfalse \\
lost' = zfalse \\
unsafeDelete' = zfalse \\
skippedOnPartial' = zfalse \\
inversion' = zfalse \\
syncTouched' = zfalse \\
syncClash' = zfalse \\
labelLost' = zfalse
\end{schema}

The name starts unowned or owned; when owned, generation~0 may hold unmarked
sync rows; a future-key generation may be present; marked sync rows may be
present even when the registry does not list the name (a lost registry).

% ===================================================================
\section{Witness tracking}
\label{sec:track}
% ===================================================================

\field{Track} is conjoined to every funnel, heal, and old-protocol step. It
latches the row-level witnesses from the before and after row maps alone, so
that no operation can forget to report a violation.

\begin{itemize}
  \item I1: a step leaves a generation in $1\upto2$ with a proper, nonempty
    part of a document \emph{and changed it}.
  \item I2: a complete funnel copy existed before and none exists after.
  \item I4: a complete funnel generation lost rows while no newer complete
    generation survives, or a writer still inserting had its rows reduced by
    someone else.
  \item I6: the newest complete funnel key after the step is smaller than
    before.
  \item Ownership: the step changed the marked sync rows or, while
    \field{g0sync} holds, the unmarked ones (\field{syncTouched}); or added a
    funnel generation while the name was owned (\field{syncClash}).
\end{itemize}

\begin{schema}{Track}
\Delta Store
\where
(\exists g : 1 \upto 2 @ 0 < rows'~g \land rows'~g < 2 \land rows'~g \neq rows~g)
  \implies partialVisible' = ztrue \\
\lnot (\exists g : 1 \upto 2 @ 0 < rows'~g \land rows'~g < 2 \land
  rows'~g \neq rows~g) \implies partialVisible' = partialVisible \\
((\exists g : 0 \upto 2 @ rows~g = 2) \land
  \lnot (\exists g : 0 \upto 2 @ rows'~g = 2)) \implies lost' = ztrue \\
\lnot ((\exists g : 0 \upto 2 @ rows~g = 2) \land
  \lnot (\exists g : 0 \upto 2 @ rows'~g = 2)) \implies lost' = lost \\
((\exists g : 0 \upto 2 @ rows~g = 2 \land rows'~g < 2 \land
    \lnot (\exists h : 0 \upto 2 @ h > g \land rows'~h = 2)) \lor
  (\exists v : WRITER @ pc~v = inserting \land pc'~v = inserting \land
    rows'~(wgen~v) < rows~(wgen~v))) \implies unsafeDelete' = ztrue \\
\lnot ((\exists g : 0 \upto 2 @ rows~g = 2 \land rows'~g < 2 \land
    \lnot (\exists h : 0 \upto 2 @ h > g \land rows'~h = 2)) \lor
  (\exists v : WRITER @ pc~v = inserting \land pc'~v = inserting \land
    rows'~(wgen~v) < rows~(wgen~v))) \implies unsafeDelete' = unsafeDelete \\
(\exists g : 0 \upto 2 @ rows~g = 2 \land
  (\forall h : g \upto 2 @ rows'~h \neq 2)) \implies inversion' = ztrue \\
\lnot (\exists g : 0 \upto 2 @ rows~g = 2 \land
  (\forall h : g \upto 2 @ rows'~h \neq 2)) \implies inversion' = inversion \\
(rows'~4 \neq rows~4 \lor (g0sync = ztrue \land rows'~0 \neq rows~0))
  \implies syncTouched' = ztrue \\
\lnot (rows'~4 \neq rows~4 \lor (g0sync = ztrue \land rows'~0 \neq rows~0))
  \implies syncTouched' = syncTouched \\
((\exists g : 1 \upto 2 @ rows'~g > rows~g) \land
  (reg = ztrue \lor rows~4 > 0)) \implies syncClash' = ztrue \\
\lnot ((\exists g : 1 \upto 2 @ rows'~g > rows~g) \land
  (reg = ztrue \lor rows~4 > 0)) \implies syncClash' = syncClash
\end{schema}

\field{KeepSkip}: the funnel steps leave \field{skippedOnPartial}, the registry
claim, and the unmarked-sync flag alone.

\begin{schema}{KeepSkip}
\Delta Store
\where
skippedOnPartial' = skippedOnPartial \\
reg' = reg \\
g0sync' = g0sync \\
fsrc' = fsrc \\
labelLost' = labelLost
\end{schema}

\field{KeepAll} is used by the sync, operator, and environment steps, which are
not under test: it leaves every witness alone.

\begin{schema}{KeepAll}
\Delta Store
\where
partialVisible' = partialVisible \\
lost' = lost \\
unsafeDelete' = unsafeDelete \\
skippedOnPartial' = skippedOnPartial \\
inversion' = inversion \\
syncTouched' = syncTouched \\
syncClash' = syncClash \\
labelLost' = labelLost
\end{schema}

% ===================================================================
\section{Operations: the proposed protocol (mode 0)}
\label{sec:ops}
% ===================================================================

A funnel writer may start a new attempt from \field{idle}, \field{done},
\field{failed}, or \field{crashed}.

\subsection{Begin, skip, refuse}

\field{BeginOverwrite} is a capture-stamped overwrite: it picks any key $k?$
that no other active attempt holds and that has no rows yet, and begins only
when the name is \emph{not owned}. \field{BeginSubmission} is a
submission-order overwrite: its key is above every complete funnel key, so it
can never be superseded. \field{BeginCreate} is a create-if-absent: no
\emph{complete} copy may exist; the future-key generation does not count as
held. \field{SkipHeld} skips a create over a complete copy.
\field{RefuseSync} is the refusal of an owned name, \field{LookupFails} the
fail-closed outcome when the registry cannot be read: both end the attempt
with no row changed.

\begin{schema}{BeginOverwriteCore}
\Delta Store \\ w? : WRITER \\ k? : 1 \upto 2
\where
mode = 0 \\
pc~w? \in \{ idle, done, failed, crashed \} \\
reg = zfalse \\
rows~4 = 0 \\
rows~k? = 0 \\
\forall v : WRITER | v \neq w? @
  (pc~v \in \{ idle, done, failed, crashed \} \lor wgen~v \neq k?) \\
pc' = pc \oplus \{ w? \mapsto staging \} \\
wgen' = wgen \oplus \{ w? \mapsto k? \} \\
seen' = seen \\
rows' = rows
\end{schema}

\begin{zed}
BeginOverwrite \defs BeginOverwriteCore \land Track \land KeepSkip
\end{zed}

\begin{schema}{BeginSubmissionCore}
BeginOverwriteCore
\where
\forall g : 0 \upto 2 @ rows~g = 2 \implies g < k?
\end{schema}

\begin{zed}
BeginSubmission \defs BeginSubmissionCore \land Track \land KeepSkip
\end{zed}

\begin{schema}{BeginCreateCore}
BeginOverwriteCore
\where
\lnot (\exists g : 0 \upto 2 @ rows~g = 2)
\end{schema}

\begin{zed}
BeginCreate \defs BeginCreateCore \land Track \land KeepSkip
\end{zed}

\begin{schema}{SkipHeldCore}
\Delta Store \\ w? : WRITER
\where
mode = 0 \\
pc~w? \in \{ idle, done, failed, crashed \} \\
\exists g : 0 \upto 2 @ rows~g = 2 \\
pc' = pc \oplus \{ w? \mapsto done \} \\
wgen' = wgen \\
seen' = seen \\
rows' = rows
\end{schema}

\begin{zed}
SkipHeld \defs SkipHeldCore \land Track \land KeepSkip
\end{zed}

A capture-stamped overwrite whose key is older than, \emph{or equal to}, a
stored \emph{complete} copy may be skipped before any work is done (an equal
stamp is the same capture again; skipping removes any tie in which a random
token could delete a stored copy). Only a complete copy counts: a damaged,
partial, or future-key one contributes no stamp. This is an optimisation and
the report, not what makes I6 hold, so \field{BeginOverwrite} stays enabled in
the same state.

\begin{schema}{SkipSupersededCore}
\Delta Store \\ w? : WRITER \\ k? : 1 \upto 2
\where
mode = 0 \\
pc~w? \in \{ idle, done, failed, crashed \} \\
\exists g : 0 \upto 2 @ g \geq k? \land rows~g = 2 \\
pc' = pc \oplus \{ w? \mapsto done \} \\
wgen' = wgen \\
seen' = seen \\
rows' = rows
\end{schema}

\begin{zed}
SkipSuperseded \defs SkipSupersededCore \land Track \land KeepSkip
\end{zed}

A submission-order write (remember, learn, inline, file) is never superseded:
its stamp is above the stored stamp. When the stored complete copy already sits
at the write ceiling (key~2), there is no higher key a write may carry and the
write \emph{fails loudly} with a refusal; it does not tie, and it does not
delete the stored copy.

\begin{schema}{RefuseCeilingCore}
\Delta Store \\ w? : WRITER
\where
mode = 0 \\
pc~w? \in \{ idle, done, failed, crashed \} \\
reg = zfalse \\
rows~4 = 0 \\
rows~2 = 2 \\
pc' = pc \oplus \{ w? \mapsto done \} \\
wgen' = wgen \\
seen' = seen \\
rows' = rows
\end{schema}

\begin{zed}
RefuseCeiling \defs RefuseCeilingCore \land Track \land KeepSkip
\end{zed}

\begin{schema}{RefuseSyncCore}
\Delta Store \\ w? : WRITER
\where
mode = 0 \\
pc~w? \in \{ idle, done, failed, crashed \} \\
(reg = ztrue \lor rows~4 > 0) \\
pc' = pc \oplus \{ w? \mapsto done \} \\
wgen' = wgen \\
seen' = seen \\
rows' = rows
\end{schema}

\begin{zed}
RefuseSync \defs RefuseSyncCore \land Track \land KeepSkip
\end{zed}

\begin{schema}{LookupFailsCore}
\Delta Store \\ w? : WRITER
\where
mode = 0 \\
pc~w? \in \{ idle, done, failed, crashed \} \\
pc' = pc \oplus \{ w? \mapsto failed \} \\
wgen' = wgen \\
seen' = seen \\
rows' = rows
\end{schema}

\begin{zed}
LookupFails \defs LookupFailsCore \land Track \land KeepSkip
\end{zed}

\subsection{Seal}

The windows are staged in the private spool; sealing completes it. Nothing is
added to \field{rows}.

\begin{schema}{SealCore}
\Delta Store \\ w? : WRITER
\where
mode = 0 \\
pc~w? = staging \\
pc' = pc \oplus \{ w? \mapsto ready \} \\
wgen' = wgen \\
seen' = seen \\
rows' = rows
\end{schema}

\begin{zed}
Seal \defs SealCore \land Track \land KeepSkip
\end{zed}

\subsection{Add}

One LanceDB commit makes the whole generation visible. Its guard is the
ownership re-check: enabled only when the name is not owned (one step here; in
the implementation the registry read and the commit run under the one
in-process write lock that sync's inserts also take).

\begin{schema}{AddCore}
\Delta Store \\ w? : WRITER
\where
mode = 0 \\
pc~w? = ready \\
reg = zfalse \\
rows~4 = 0 \\
rows' = rows \oplus \{ wgen~w? \mapsto 2 \} \\
pc' = pc \oplus \{ w? \mapsto healing \} \\
wgen' = wgen \\
seen' = seen
\end{schema}

\begin{zed}
Add \defs AddCore \land Track \land KeepSkip
\end{zed}

\subsection{Heal: ownership check, observe, then delete}

A heal \emph{starts} only on a name that is not owned (the check is the
registry lookup; the writer's own heal, the held-check heal, and the sweep all
begin here), observes the newest complete funnel generation, then deletes every
funnel generation \emph{older} than it, every \emph{damaged} generation of any
key (partial, or the future key), and nothing else. It never deletes the newest
complete generation, never deletes when none was complete, never touches marked
sync rows, and re-verifies the observed generation before the delete.

\begin{schema}{BeginHealCore}
\Delta Store \\ w? : WRITER
\where
mode = 0 \\
pc~w? \in \{ idle, done, failed, crashed \} \\
reg = zfalse \\
rows~4 = 0 \\
pc' = pc \oplus \{ w? \mapsto healing \} \\
wgen' = wgen \\
seen' = seen \\
rows' = rows
\end{schema}

\begin{zed}
BeginHeal \defs BeginHealCore \land Track \land KeepSkip
\end{zed}

\begin{schema}{ObserveCore}
\Delta Store \\ w? : WRITER
\where
mode = 0 \\
pc~w? = healing \\
pc' = pc \oplus \{ w? \mapsto observed \} \\
(\exists g : 0 \upto 2 @ rows~g = 2) \implies
  (seen'~w? \leq 2 \land rows~(seen'~w?) = 2 \land
   (\forall h : 0 \upto 2 @ rows~h = 2 \implies h \leq seen'~w?)) \\
\lnot (\exists g : 0 \upto 2 @ rows~g = 2) \implies seen'~w? = 5 \\
\forall v : WRITER | v \neq w? @ seen'~v = seen~v \\
wgen' = wgen \\
rows' = rows
\end{schema}

\begin{zed}
Observe \defs ObserveCore \land Track \land KeepSkip
\end{zed}

\begin{schema}{DeleteOlderCore}
\Delta Store \\ w? : WRITER
\where
mode = 0 \\
pc~w? = observed \\
seen~w? = 5 \implies rows' = rows \\
(seen~w? \leq 2 \land rows~(seen~w?) \neq 2) \implies rows' = rows \\
(seen~w? \leq 2 \land rows~(seen~w?) = 2) \implies
  ((\forall g : 0 \upto 2 @
    ((g < seen~w? \lor (g \geq 1 \land 0 < rows~g \land rows~g < 2))
       \implies rows'~g = 0) \land
    (\lnot (g < seen~w? \lor (g \geq 1 \land 0 < rows~g \land rows~g < 2))
       \implies rows'~g = rows~g)) \land
   rows'~3 = 0 \land rows'~4 = rows~4) \\
pc' = pc \oplus \{ w? \mapsto done \} \\
wgen' = wgen \\
seen' = seen
\end{schema}

\begin{zed}
DeleteOlder \defs DeleteOlderCore \land Track \land KeepSkip
\end{zed}

\subsection{Fail and crash}

A failure at any step before \field{Add} touches no row. A failure or crash
after \field{Add} leaves the committed generation alone; older generations stay
until the next heal.

\begin{schema}{FailCore}
\Delta Store \\ w? : WRITER
\where
mode = 0 \\
pc~w? \in \{ staging, ready, healing, observed \} \\
pc' = pc \oplus \{ w? \mapsto failed \} \\
wgen' = wgen \\
seen' = seen \\
rows' = rows
\end{schema}

\begin{zed}
Fail \defs FailCore \land Track \land KeepSkip
\end{zed}

\begin{schema}{CrashCore}
\Delta Store \\ w? : WRITER
\where
pc~w? \in \{ staging, ready, healing, observed, inserting \} \\
pc' = pc \oplus \{ w? \mapsto crashed \} \\
wgen' = wgen \\
seen' = seen \\
rows' = rows
\end{schema}

\begin{zed}
Crash \defs CrashCore \land Track \land KeepSkip
\end{zed}

% ===================================================================
\section{Sync, operator, and environment}
\label{sec:other}
% ===================================================================

Enabled in both modes; none is under test.

\subsection{Sync}

Sync writes only names its registry covers, and a name enters the registry's
\emph{files} table when sync first writes it. Its \emph{first} write of a name
that holds funnel rows is a \emph{takeover} only when those rows are
file-sourced from the same path (\field{fsrc}): it then deletes them by name
(the registered directory is authoritative for its files) --- funnel
generations $1\upto3$, and generation~0 unless it is sync's own unmarked rows.
When the funnel rows are anything else (a label, a capture, a URL page) sync
\emph{skips the file} and records a conflict (\field{SyncConflict}, which
changes nothing the model tracks); the funnel copy stays and the name stays
unowned. Sync's rows are marked, window by window, and a partial sync-owned
document is legitimate: the registry owns completeness.

\begin{schema}{SyncInsertCore}
\Delta Store
\where
rows~4 < 2 \\
((\exists g : 1 \upto 3 @ rows~g > 0) \lor
  (g0sync = zfalse \land rows~0 > 0)) \implies fsrc = ztrue \\
rows'~4 = rows~4 + 1 \\
rows'~1 = 0 \\
rows'~2 = 0 \\
rows'~3 = 0 \\
g0sync = ztrue \implies rows'~0 = rows~0 \\
g0sync = zfalse \implies rows'~0 = 0 \\
pc' = pc \\
wgen' = wgen \\
seen' = seen \\
reg' = ztrue \\
g0sync' = g0sync \\
fsrc' = fsrc \\
((\exists g : 1 \upto 3 @ rows~g > 0) \lor
  (g0sync = zfalse \land rows~0 > 0)) \land fsrc = zfalse
  \implies labelLost' = ztrue \\
\lnot (((\exists g : 1 \upto 3 @ rows~g > 0) \lor
  (g0sync = zfalse \land rows~0 > 0)) \land fsrc = zfalse)
  \implies labelLost' = labelLost \\
partialVisible' = partialVisible \\
lost' = lost \\
unsafeDelete' = unsafeDelete \\
skippedOnPartial' = skippedOnPartial \\
inversion' = inversion \\
syncTouched' = syncTouched \\
syncClash' = syncClash
\end{schema}

\begin{zed}
SyncInsert \defs SyncInsertCore
\end{zed}

\begin{schema}{SyncConflictCore}
\Delta Store
\where
reg = zfalse \\
((\exists g : 1 \upto 3 @ rows~g > 0) \lor
  (g0sync = zfalse \land rows~0 > 0)) \\
fsrc = zfalse \\
rows' = rows \\
pc' = pc \\
wgen' = wgen \\
seen' = seen \\
reg' = reg \\
g0sync' = g0sync \\
fsrc' = fsrc
\end{schema}

\begin{zed}
SyncConflict \defs SyncConflictCore \land KeepAll
\end{zed}

\begin{schema}{SyncDeleteCore}
\Delta Store
\where
rows~4 > 0 \\
rows'~4 = 0 \\
g0sync = ztrue \implies rows'~0 = 0 \\
g0sync = zfalse \implies rows'~0 = rows~0 \\
rows'~1 = rows~1 \\
rows'~2 = rows~2 \\
rows'~3 = rows~3 \\
pc' = pc \\
wgen' = wgen \\
seen' = seen \\
reg' = reg \\
g0sync' = zfalse \\
fsrc' = fsrc
\end{schema}

\begin{zed}
SyncDelete \defs SyncDeleteCore \land KeepAll
\end{zed}

\subsection{Operator}

\field{Purge} deregisters a directory and deletes its data: the name leaves the
registry. With the data kept, deregistering changes nothing the model tracks,
because the name stays in the retained set.

\begin{schema}{PurgeCore}
\Delta Store
\where
reg = ztrue \\
reg' = zfalse \\
g0sync' = zfalse \\
rows'~4 = 0 \\
g0sync = ztrue \implies rows'~0 = 0 \\
g0sync = zfalse \implies rows'~0 = rows~0 \\
rows'~1 = rows~1 \\
rows'~2 = rows~2 \\
rows'~3 = rows~3 \\
pc' = pc \\
wgen' = wgen \\
seen' = seen \\
fsrc' = fsrc
\end{schema}

\begin{zed}
Purge \defs PurgeCore \land KeepAll
\end{zed}

\subsection{Environment}

One row of a funnel generation disappears: an external delete, a mixed-version
writer, a delete that half-failed. A complete generation becomes damaged; the
heal must repair it.

\begin{schema}{ExternalDeleteCore}
\Delta Store \\ g? : 0 \upto 2
\where
mode = 0 \\
rows~g? > 0 \\
rows' = rows \oplus \{ g? \mapsto rows~g? - 1 \} \\
reg' = reg \\
g0sync' = g0sync \\
fsrc' = fsrc \\
pc' = pc \\
wgen' = wgen \\
seen' = seen
\end{schema}

\begin{zed}
ExternalDelete \defs ExternalDeleteCore \land KeepAll
\end{zed}

% ===================================================================
\section{Operations: the previous protocol (mode 1)}
\label{sec:old}
% ===================================================================

These operations reproduce the funnel as it was: overwrite deletes every
prior row of the name whatever its owner or key, windows insert straight into
the real name, a failed insert deletes the name's rows, create-if-absent
treats any row under the name as held, and nothing compares capture times,
owners, or the registry. They are enabled only when \field{mode}~$=1$.

\begin{schema}{OldBeginOverwriteCore}
\Delta Store \\ w? : WRITER \\ k? : 1 \upto 2
\where
mode = 1 \\
pc~w? \in \{ idle, done, failed, crashed \} \\
\forall v : WRITER | v \neq w? @
  (pc~v \in \{ idle, done, failed, crashed \} \lor wgen~v \neq k?) \\
\forall g : 0 \upto 4 @ rows'~g = 0 \\
pc' = pc \oplus \{ w? \mapsto inserting \} \\
wgen' = wgen \oplus \{ w? \mapsto k? \} \\
seen' = seen \\
skippedOnPartial' = skippedOnPartial \\
reg' = reg \\
g0sync' = zfalse \\
fsrc' = fsrc \\
labelLost' = labelLost
\end{schema}

\begin{zed}
OldBeginOverwrite \defs OldBeginOverwriteCore \land Track
\end{zed}

\begin{schema}{OldBeginCreateCore}
\Delta Store \\ w? : WRITER \\ k? : 1 \upto 2
\where
mode = 1 \\
pc~w? \in \{ idle, done, failed, crashed \} \\
\forall v : WRITER | v \neq w? @
  (pc~v \in \{ idle, done, failed, crashed \} \lor wgen~v \neq k?) \\
\forall g : 0 \upto 4 @ rows~g = 0 \\
rows' = rows \\
pc' = pc \oplus \{ w? \mapsto inserting \} \\
wgen' = wgen \oplus \{ w? \mapsto k? \} \\
seen' = seen
\end{schema}

\begin{zed}
OldBeginCreate \defs OldBeginCreateCore \land Track \land KeepSkip
\end{zed}

The old held check: any row under the name means ``already stored''. A
partial copy is skipped forever; \field{skippedOnPartial} latches.

\begin{schema}{OldSkipHeldCore}
\Delta Store \\ w? : WRITER
\where
mode = 1 \\
pc~w? \in \{ idle, done, failed, crashed \} \\
\exists g : 0 \upto 4 @ rows~g > 0 \\
pc' = pc \oplus \{ w? \mapsto done \} \\
rows' = rows \\
wgen' = wgen \\
seen' = seen \\
reg' = reg \\
g0sync' = g0sync \\
fsrc' = fsrc \\
labelLost' = labelLost \\
(\lnot (\exists g : 0 \upto 2 @ rows~g = 2)) \implies skippedOnPartial' = ztrue \\
(\exists g : 0 \upto 2 @ rows~g = 2) \implies
  skippedOnPartial' = skippedOnPartial
\end{schema}

\begin{zed}
OldSkipHeld \defs OldSkipHeldCore \land Track
\end{zed}

\begin{schema}{OldInsertCore}
\Delta Store \\ w? : WRITER
\where
mode = 1 \\
pc~w? = inserting \\
rows~(wgen~w?) < 2 \\
rows' = rows \oplus \{ wgen~w? \mapsto rows~(wgen~w?) + 1 \} \\
pc' = pc \\
wgen' = wgen \\
seen' = seen
\end{schema}

\begin{zed}
OldInsert \defs OldInsertCore \land Track \land KeepSkip
\end{zed}

\begin{schema}{OldFinishCore}
\Delta Store \\ w? : WRITER
\where
mode = 1 \\
pc~w? = inserting \\
rows~(wgen~w?) = 2 \\
pc' = pc \oplus \{ w? \mapsto done \} \\
rows' = rows \\
wgen' = wgen \\
seen' = seen
\end{schema}

\begin{zed}
OldFinish \defs OldFinishCore \land Track \land KeepSkip
\end{zed}

The old failure cleanup: delete by name, whatever owner or generation the
rows belong to.

\begin{schema}{OldFailCore}
\Delta Store \\ w? : WRITER
\where
mode = 1 \\
pc~w? = inserting \\
\forall g : 0 \upto 4 @ rows'~g = 0 \\
pc' = pc \oplus \{ w? \mapsto failed \} \\
wgen' = wgen \\
seen' = seen \\
skippedOnPartial' = skippedOnPartial \\
reg' = reg \\
g0sync' = zfalse \\
fsrc' = fsrc \\
labelLost' = labelLost
\end{schema}

\begin{zed}
OldFail \defs OldFailCore \land Track
\end{zed}

The previous sync took a name over unconditionally (this was the earlier
ruling, before the file-sourced rule): it is the mode-1 operation that reaches
\field{labelLost}.

\begin{schema}{OldSyncTakeoverCore}
\Delta Store
\where
mode = 1 \\
rows~4 < 2 \\
((\exists g : 1 \upto 3 @ rows~g > 0) \lor
  (g0sync = zfalse \land rows~0 > 0)) \\
rows'~4 = rows~4 + 1 \\
rows'~1 = 0 \\
rows'~2 = 0 \\
rows'~3 = 0 \\
g0sync = ztrue \implies rows'~0 = rows~0 \\
g0sync = zfalse \implies rows'~0 = 0 \\
pc' = pc \\
wgen' = wgen \\
seen' = seen \\
reg' = ztrue \\
g0sync' = g0sync \\
fsrc' = fsrc \\
fsrc = zfalse \implies labelLost' = ztrue \\
fsrc = ztrue \implies labelLost' = labelLost \\
partialVisible' = partialVisible \\
lost' = lost \\
unsafeDelete' = unsafeDelete \\
skippedOnPartial' = skippedOnPartial \\
inversion' = inversion \\
syncTouched' = syncTouched \\
syncClash' = syncClash
\end{schema}

\begin{zed}
OldSyncTakeover \defs OldSyncTakeoverCore
\end{zed}

% ===================================================================
\section{Properties checked}
\label{sec:props}
% ===================================================================

Each property is a goal that ProB must \emph{not} reach in mode~0 and
\emph{must} reach in mode~1 (where applicable). All claims are scoped to this
model (\S\ref{sec:nongoals}).

\begin{center}
\begin{tabular}{@{}lll@{}}
\toprule
Property & Goal (unreachable in mode 0) & Meaning \\
\midrule
I1 & \field{partialVisible = ztrue} & a protocol step produced a partial copy \\
I2 & \field{lost = ztrue} & a complete copy existed and none survives \\
I4 & \field{unsafeDelete = ztrue} & cleanup removed a copy it did not supersede \\
create & \field{skippedOnPartial = ztrue} & held check trusted a partial copy \\
I6 & \field{inversion = ztrue} & the newest complete key went backwards \\
own & \field{syncTouched = ztrue} & funnel or heal changed sync rows, marked or not \\
own & \field{syncClash = ztrue} & funnel added a generation to an owned name \\
own & \field{labelLost = ztrue} & sync's takeover deleted a label, capture, or URL copy \\
\bottomrule
\end{tabular}
\end{center}

I3 (replace converges) has two parts. The transient \emph{duplicate} (two
complete generations) is reachable and is reported as such. Convergence is the
CTL formula \field{AG EF} that from every reachable state a quiescent state (no
active writer) with at most one complete generation, and, when one is
complete, no damaged funnel generation in $1\upto2$ and no future-key
generation, is reachable. (Without a complete generation a damaged one is
deliberately kept for the doctor; the model has two keys, so a store in which
both are damaged and none is complete has no free key. The real key space has
no such dead end.) This is a \emph{possibility} claim: it says no reachable
state is a dead end for repair, not that every execution reaches a clean state.

\end{document}
